\begin{tikzpicture}[>=Latex,x=1cm,y=1cm]
  % Mode-specific bounds keep the calculations optically centered. The extra
  % margin prevents the labels and the first/last calculation lines from being
  % cropped in the cached standalone PDF.
  \IfClozeTF
    {\path[use as bounding box] (-0.9,-1.7) rectangle (18.6,9.7);}
    {\path[use as bounding box] (-0.9,-1.7) rectangle (16.2,9.7);}

  \node[anchor=west,inner sep=0pt] (rowone) at (3.35,8) {$
    \begin{aligned}
      (\vec a\times\vec b)\cdot\vec c
      &=\begin{ClozeContinuation}[21]{3}
        \left[
          \begin{pmatrix}3\\-1\\2\end{pmatrix}
          \times
          \begin{pmatrix}4\\2\\1\end{pmatrix}
        \right]\cdot\begin{pmatrix}0\\2\\-5\end{pmatrix}
      \end{ClozeContinuation}\\[0.4em]
      &=\begin{ClozeContinuation}[21]{3}
        \begin{pmatrix}-5\\5\\10\end{pmatrix}
          \cdot\begin{pmatrix}0\\2\\-5\end{pmatrix}
          =0+10-50=-40
      \end{ClozeContinuation}
    \end{aligned}$};

  \node[anchor=west,inner sep=0pt] (rowtwo) at (3.35,4) {$
    \begin{aligned}
      (\vec b\times\vec c)\cdot\vec a
      &=\begin{ClozeContinuation}[21]{3}
        \left[
          \begin{pmatrix}4\\2\\1\end{pmatrix}
          \times
          \begin{pmatrix}0\\2\\-5\end{pmatrix}
        \right]\cdot\begin{pmatrix}3\\-1\\2\end{pmatrix}
      \end{ClozeContinuation}\\[0.4em]
      &=\begin{ClozeContinuation}[21]{3}
        \begin{pmatrix}-12\\20\\8\end{pmatrix}
          \cdot\begin{pmatrix}3\\-1\\2\end{pmatrix}
          =-36-20+16=-40
      \end{ClozeContinuation}
    \end{aligned}$};

  \node[anchor=west,inner sep=0pt] (rowthree) at (3.35,0) {$
    \begin{aligned}
      (\vec a\times\vec c)\cdot\vec b
      &=\begin{ClozeContinuation}[21]{3}
        \left[
          \begin{pmatrix}3\\-1\\2\end{pmatrix}
          \times
          \begin{pmatrix}0\\2\\-5\end{pmatrix}
        \right]\cdot\begin{pmatrix}4\\2\\1\end{pmatrix}
      \end{ClozeContinuation}\\[0.4em]
      &=\begin{ClozeContinuation}[21]{3}
        \begin{pmatrix}1\\15\\6\end{pmatrix}
          \cdot\begin{pmatrix}4\\2\\1\end{pmatrix}
          =4+30+6=40
      \end{ClozeContinuation}
    \end{aligned}$};

  \coordinate (arrowone) at ([xshift=-3mm,yshift=10mm]rowone.west);
  \coordinate (arrowtwoabove) at ([xshift=-3mm,yshift=11.5mm]rowtwo.west);
  \coordinate (arrowtwobelow) at ([xshift=-3mm,yshift=8.5mm]rowtwo.west);
  \coordinate (arrowthree) at ([xshift=-3mm,yshift=10mm]rowthree.west);

  \draw[<->,THMBlue,line width=1.1pt]
    (arrowone) .. controls +(-1.0,0) and +(-1.0,0) .. (arrowtwoabove)
    node[midway,left=0.3cm,align=center,text=THMBlue]
      {zyklische\\Vertauschung};

  \draw[<->,olive!80!black,line width=1.1pt]
    (arrowtwobelow) .. controls +(-1.0,0) and +(-1.0,0) .. (arrowthree)
    node[midway,left=0.3cm,align=center,text=olive!80!black]
      {paarweise\\Vertauschung};
\end{tikzpicture}
